% !TEX root = ../../thesis.tex
\chapter{Thực nghiệm và đánh giá}

Chương này kiểm chứng thực nghiệm hệ thống Fiber đã được xây dựng trong Chương~3 thông qua khung đánh giá ba trục trên tập dữ liệu gồm 154 bài báo tiếng Việt cân bằng theo năm nhóm chủ đề. Mục~4.1 trình bày thiết lập thực nghiệm bao gồm tập dữ liệu, các mô hình so sánh và hệ đo lường; Mục~4.2 báo cáo kết quả Trục~A về khả năng giữ lại nội dung; Mục~4.3 trình bày kết quả Trục~B về chất lượng và mức độ ưu tiên, bao gồm phép so sánh mang tính quyết định giữa bản dung hợp MoA và GPT-4o đơn lẻ; Mục~4.4 phân tích kết quả Trục~C từ khảo sát 17 người đánh giá độc lập; Mục~4.5 thực hiện phân tích đa trục để đối chiếu và tam giác hóa các kết quả; và Mục~4.6 thảo luận về các hạn chế phương pháp luận.

\section{Thiết lập thực nghiệm}

\subsection{Tập dữ liệu}

Quá trình đánh giá hiệu năng của các mô hình được thực hiện trên tập dữ liệu gồm $154$ bài báo tiếng Việt, được thu thập trực tiếp từ báo điện tử \textit{tienphong.vn}. Để đảm bảo tính toàn diện và khách quan, tập dữ liệu được thiết kế cân bằng cấu trúc theo chủ đề với tỷ lệ phân bổ xấp xỉ $30$ bài mỗi nhóm, trải dài trên năm danh mục nội dung cốt lõi bao gồm: \textit{Thời sự}, \textit{Pháp luật}, \textit{Kinh tế}, \textit{Giáo dục} và \textit{Văn hóa}. 

Việc thiết kế một tập dữ liệu cân bằng theo phân vùng chủ đề (được chi tiết hóa trong Bảng~\ref{tab:dataset_distribution}) đóng vai trò quan trọng trong việc kiểm soát các biến nhiễu do phong cách viết đặc thù, thuật ngữ chuyên ngành hoặc độ dài trung bình của từng thể loại báo chí. Đồng thời, cấu trúc này thiết lập cơ sở dữ liệu nền tảng phục vụ cho các phân tích sâu về sự tương tác và biến thiên hiệu năng của mô hình theo từng nhóm chủ đề được trình bày chi tiết tại Mục~4.4.3.

\begin{table}[htbp]
\centering
\caption{Phân bố số lượng bài báo trong tập dữ liệu thực nghiệm theo danh mục}
\label{tab:dataset_distribution}
% Đổi từ tabular sang tabularx, set width = \textwidth và dùng cột X
\begin{tabularx}{\textwidth}{l c X}
\hline
Chủ đề & Số lượng bài báo & Đặc trưng nội dung \\ \hline
Thời sự         & 31                        & Tin tức chính trị, xã hội, sự kiện nóng. \\
Pháp luật       & 31                        & Văn bản pháp quy, án lệ, an ninh trật tự. \\
Kinh tế         & 30                        & Tài chính, thị trường, kinh doanh vĩ mô. \\
Giáo dục        & 31                        & Tuyển sinh, chính sách giáo dục, học đường. \\
Văn hóa         & 31                        & Nghệ thuật, đời sống, giá trị truyền thống. \\ \hline
Tổng cộng & 154        &  \\ \hline
\end{tabularx}
\end{table}

\subsection{Các mô hình so sánh}

Năm mô hình xử lý tóm tắt văn bản được đưa vào đánh giá và so sánh trên từng bài báo:

\begin{itemize}
  \item MoA FUSION --- gồm ba mô hình GPT-4o-mini, Claude Haiku 4.5, Gemini 2.5 Flash là tác nhân đề xuất kết hợp với tác nhân tổng hợp GPT-4o. 
  \item GPT-4o --- mô hình cơ sở mạnh nhất, sử dụng mô hình \texttt{gpt-4o-2024-08-06}. Phương án này được sử dụng trong phép so sánh mang tính quyết định của khóa luận.
  \item Claude Haiku 4.5 --- mô hình cơ sở thuộc phân cấp tác nhân đề xuất (nhà cung cấp Anthropic).
  \item Gemini 2.5 Flash --- mô hình cơ sở thuộc phân cấp tác nhân đề xuất (nhà cung cấp Google).
  \item GPT-4o-mini --- mô hình cơ sở thuộc phân cấp tác nhân đề xuất (nhà cung cấp OpenAI).
\end{itemize}

\subsection{Các hệ đo lường}

\begin{itemize}
  \item Trục A (khả năng giữ lại nội dung): Bao gồm các hệ số ROUGE-1, ROUGE-L, BLEU, BERTScore (với mô hình nền PhoBERT) và tỷ lệ nén. Các hệ số này được tính toán so với bài viết nguồn.
  \item Trục B (chất lượng và mức độ ưu tiên): Bao gồm điểm số đánh giá theo thang điểm 1--5, tỷ lệ thắng đấu cặp kèm theo giá trị $p$ của kiểm định dấu và kiểm soát phân nhóm theo độ dài, tỷ lệ chấp thuận logic phục vụ xác định tính xác thực. Mô hình LLM-Judge được sử dụng là \texttt{gpt-4o-mini-2024-07-18}.
  \item Trục C (đánh giá của con người): Phương pháp xếp hạng mù $K = 3$; tính toán thứ hạng trung bình của từng phương án, tỷ lệ xếp hạng nhất, tỷ lệ thắng đấu cặp, và đo lường độ đồng thuận giữa các người đánh giá.
\end{itemize} 

\section{Kết quả Trục A --- Khả năng giữ lại nội dung}

Kết quả thực nghiệm đo được của trục A được biểu thị trong bảng~\ref{tab:axisA} cho thấy báo cáo các hệ số overlap từ vựng được tính trung bình trên 154 lượt chạy dung hợp và các lượt chạy mô hình đơn lẻ tương đương làm cơ sở đối chiếu.

\begin{table}[H]
  \centering
  \small
  \setlength{\tabcolsep}{4pt}
  \caption{Trục A --- các hệ số overlap so với bài viết nguồn (điểm số càng cao thể hiện bản tóm tắt càng bám sát văn phong từ vựng của bài gốc).}
  \label{tab:axisA}
  \begin{tabularx}{\linewidth}{Xrrrrrr}
    \toprule
    Mô hình & $n$ & ROUGE-1 & ROUGE-L & BLEU & BERTScore & Tỷ lệ nén\% \\
    \midrule
    \textbf{MoA FUSION}          & \textbf{154} & \textbf{0,4240} & \textbf{0,3300} & \textbf{0,1349} & \textbf{0,6759} & \textbf{42,6} \\
    gemini-2.5-flash             & 100 & 0,3503 & 0,2677 & 0,0894 & 0,6538 & 33,7 \\
    claude-haiku-4.5             & 100 & 0,3502 & 0,2830 & 0,0871 & 0,6487 & 34,7 \\
    gpt-4o                       & 50  & 0,3621 & 0,2567 & 0,0688 & 0,6246 & 37,6 \\
    gpt-4o-mini                  & 100 & 0,2454 & 0,1849 & 0,0290 & 0,6217 & 24,9 \\
    \bottomrule
  \end{tabularx}
\end{table}

Mô hình kết hợp MoA FUSION chiến thắng trên tất cả các hệ số đo lường mức độ giữ lại nội dung --- đạt điểm ROUGE-1 $0,4240$ so với $0,3621$ của GPT-4o đơn lẻ, và BERTScore $0,6759$ so với $0,6246$. Điều này \textit{trái ngược} với những dự đoán từ nghiên cứu thử nghiệm ban đầu của khóa luận (khi đó khoá luận giả định rằng kết quả dung hợp MoA sẽ có điểm overlap thấp hơn do cơ chế tổng hợp kiểu biên tập sẽ diễn giải lại thông tin một cách mạnh mẽ hơn).

Dự đoán trong giai đoạn thử nghiệm ban đầu được đưa ra dựa trên một câu lệnh tổng hợp nguồn nghiêm ngặt (được thử nghiệm trước 08-05-2026). Dưới tác động của câu lệnh ràng buộc nghiêm ngặt đó, mô hình tổng hợp đã thu gọn các bản thảo thành một đoạn văn biên tập rất ngắn, trực tiếp kéo điểm overlap đi xuống. Sau khi tối ưu hóa câu lệnh tổng hợp, mô hình tổng hợp giờ đây vừa giữ lại được khung lập luận trung thực với bài gốc, vừa diễn giải các dữ kiện chính bằng các cụm từ rất sát với nguồn, khiến điểm overlap tăng vượt trội hơn bất kỳ bản thảo đơn lẻ nào.

Kết quả sau khi tối ưu hóa câu lệnh mang ý nghĩa thực tiễn quan trọng: ngay cả trên Trục A, phương pháp dung hợp vẫn cho thấy hiệu quả tương đương hoặc vượt trội so với mô hình đơn lẻ mạnh nhất. Hệ đo lường trùng lặp so với bài gốc không còn là một điểm yếu; trái lại, nó trở thành trục thứ ba đồng thuận chặt chẽ với Trục B và Trục C, giúp củng cố vững chắc hơn luận điểm khoa học ủng hộ phương pháp dung hợp đa tác nhân.

\section{Kết quả Trục B --- Chất lượng và sự ưu tiên}

\subsection{Đánh giá theo Rubric (FLASK cải tiến, 1--5)}

Bảng~\ref{tab:axisB1} báo cáo điểm số rubric trung bình theo từng chiều đánh giá trên 154 lượt chạy dung hợp và các lượt chạy mô hình đơn lẻ tương ứng.

\begin{table}[H]
  \centering
  \caption{Trục B.1 --- điểm số rubric đánh giá bằng LLM-giám khảo (1--5). Sử dụng \texttt{judge\_model = gpt-4o-mini-2024-07-18}.}
  \label{tab:axisB1}
  \begin{tabularx}{\linewidth}{Xrrrrrr}
    \toprule
    Mô hình & $n$ & Faithfulness & Coverage & Fluency & Conciseness & Overall \\
    \midrule
    \textbf{MoA FUSION}          & \textbf{154} & \textbf{5,00} & \textbf{4,94} & \textbf{4,98} & \textbf{4,88} & \textbf{4,96} \\
    claude-haiku-4.5             & 100 & 5,00 & 4,78 & 4,97 & 4,92 & 4,91 \\
    gpt-4o                       & 50  & 5,00 & 4,78 & 4,98 & 4,84 & 4,88 \\
    gemini-2.5-flash             & 100 & 4,96 & 4,76 & 5,00 & 4,95 & 4,86 \\
    gpt-4o-mini                  & 100 & 4,98 & 4,34 & 4,97 & 4,84 & 4,61 \\
    \bottomrule
  \end{tabularx}
\end{table}

Mô hình kết hợp MoA FUSION vượt trội hơn mọi mô hình cơ sở ở điểm số đánh giá Tổng thể toàn diện ($4,96$ so với $4,91$ của mô hình đứng sau). Yếu tố đóng góp khác biệt lớn nhất là tiêu chí Độ bao phủ (Coverage) ($4,94$ so với $4,78$ của GPT-4o đơn lẻ, tạo ra khoảng cách $+0,16$); kết quả này hoàn toàn nhất quán với cơ chế ``hợp các bản thảo'' được thảo luận trong Mục~4.2.2. Điểm số tiêu chí Tính trung thực (Faithfulness) gần như bão hòa đối với tất cả các phương pháp ($\geq 4,96$), phản ánh rằng mô hình kết hợp và các mô hình ứng viên đều là những mô hình mạnh mẽ, và hiện tượng mâu thuẫn logic trực tiếp là rất hiếm gặp.

\subsection{Kết quả dung hợp so với Bản thảo đơn lẻ tốt nhất}

Đây là phép so sánh đối chiếu đấu cặp chính trên toàn bộ tập dữ liệu gồm $154$ bài báo. Mô hình LLM-Judge sẽ đánh giá bản tóm tắt dung hợp MoA đối đầu với bản thảo đề xuất đơn lẻ mạnh nhất ứng với từng bài báo.

\begin{table}[H]
  \centering
  \caption{Trục B.2a --- phán quyết đấu cặp, bản dung hợp (A) đấu với bản thảo tốt nhất (B) trên toàn bộ tập dữ liệu. Các kết quả hòa được tính vào tổng số lượng.}
  \label{tab:axisB2a}
  \begin{tabularx}{\linewidth}{rCCCC}
    \toprule
    $n$ & Bản dung hợp thắng & Bản thảo tốt nhất thắng & Hòa & Tỷ lệ thắng dung hợp (không tính hòa) \\
    \midrule
    153 & 93 & 42 & 18 & \textbf{68,9\%} \\
    \bottomrule
  \end{tabularx}
\end{table}

Trong số $153$ phán quyết phân định rõ ràng, bản tóm tắt dung hợp giành chiến thắng $93$ lần, bản thảo đơn lẻ tốt nhất thắng $42$ lần, và có $18$ trường hợp hòa. Nếu loại trừ các trường hợp hòa, bản tóm tắt dung hợp được ưa thích hơn bản thảo tốt nhất trong gần $69\%$ các lượt so sánh trực tiếp (chiếm $60,8\%$ tổng số lượt so sánh). Mẫu thử nghiệm quy mô lớn này xác nhận rằng việc cung cấp nhiều bản thảo làm ngữ cảnh cho mô hình tổng hợp giúp tạo ra bản tóm tắt chất lượng vượt trội hơn hẳn so với việc chỉ sử dụng một mô hình đề xuất đơn lẻ.

\subsection{Kết quả dung hợp so với từng mô hình đề xuất}

Nhằm đối chiếu trực tiếp với kết quả nghiên cứu của Wang et al.\ \cite{Wang2024} (Hình 4a), Bảng~\ref{tab:axisB2b} báo cáo tỷ lệ chiến thắng của phương pháp dung hợp khi đối đầu với từng mô hình đề xuất riêng lẻ.

\begin{table}[H]
  \centering
  \caption{Trục B.2b --- bản dung hợp đấu với từng mô hình đề xuất. Kiểm định dấu (sign-test) được thực hiện hai chiều; loại trừ các trường hợp hòa.}
  \label{tab:axisB2b}
  \begin{tabularx}{\linewidth}{X r c c c C C}
    \toprule
    Mô hình đề xuất & $n$ & Dung hợp thắng & Bản thảo thắng & Hòa & Tỷ lệ thắng dung hợp & Giá trị $p$ Sign-test \\
    \midrule
    gpt-4o-mini       & 50 & 49 & 1 & 0 & 98,0\% & 0,0000 \\
    claude-haiku-4.5  & 50 & 41 & 8 & 1 & 83,7\% & 0,0000 \\
    gemini-2.5-flash  & 50 & 32 & 14 & 4 & 69,6\% & 0,0114 \\
    \bottomrule
  \end{tabularx}
\end{table}

Tỷ lệ chiến thắng thô của bản dung hợp lần lượt là $98,0\%$, $83,7\%$, $69,6\%$; sau khi áp dụng kiểm soát phân nhóm theo độ dài (length-bucketed), tỷ lệ này đạt $97,9\%$, $84,3\%$, $72,0\%$. Điểm số phân nhóm biến động rất nhỏ (trong khoảng $\pm 2,4$ điểm phần trăm so với tỷ lệ thô), xác nhận rằng thiên kiến độ dài của LLM-Judge không phải là yếu tố ảnh hưởng kết quả này.

Bản tóm tắt dung hợp vượt qua từng mô hình đề xuất riêng lẻ với ý nghĩa thống kê rõ rệt. Khoảng cách này lớn nhất khi đối đầu với mô hình yếu nhất (GPT-4o-mini, $98\%$) và thu hẹp nhất khi gặp mô hình mạnh nhất (Gemini 2.5 Flash, $70\%$). Hiện tượng này hoàn toàn khớp với đặc tính kỹ thuật được công bố trong nghiên cứu MoA gốc \cite{Wang2024}: cơ chế dung hợp tạo ra chênh lệch lớn nhất khi các mô hình đề xuất có năng lực hạn chế hơn, và sự hiện diện của một mô hình đề xuất mạnh sẽ thu hẹp nhưng không thể xóa bỏ khoảng cách năng lực này.

\subsection{Kết quả dung hợp so với GPT-4o}

Đây là phép so sánh đối chiếu có kiểm soát mô hình mang tính quyết định nhất trong khóa luận này. Bằng cách cố định mô hình tạo ra kết quả cuối cùng (cả hai bản tóm tắt đối đầu đều là kết quả đầu ra của GPT-4o), khoá luận đã có thể cô lập hoàn toàn tác động của quy trình dung hợp (cung cấp ba bản thảo đề xuất làm ngữ cảnh đầu vào cho mô hình tổng hợp) khỏi tác động của năng lực bản thân mô hình.

\begin{table}[H]
  \centering
  \caption{Trục B.2c --- phán quyết đấu cặp, bản dung hợp (A) so sánh với GPT-4o chạy đơn lẻ (B). Kiểm định dấu thực hiện hai chiều; loại trừ các trường hợp hòa.}
  \label{tab:axisB2c}
  \begin{tabularx}{\linewidth}{r c c c C C}
    \toprule
    $n$ & Dung hợp thắng & Đơn lẻ thắng & Hòa & Tỷ lệ thắng dung hợp & Giá trị $p$ Sign-test \\
    \midrule
    48 & 37 & 11 & 0 & \textbf{77,1\%} & \textbf{0,0002} \\
    \bottomrule
  \end{tabularx}
\end{table}

Trong số $48$ kết quả so sánh phân định rõ ràng, bản dung hợp giành chiến thắng $37$ lần, mô hình GPT-4o chạy đơn lẻ thắng $11$ lần, không ghi nhận trường hợp hòa nào. Dưới giả thuyết không $H_0: P(\text{dung hợp thắng}) = 0,5$, xác suất xuất hiện từ $37$ chiến thắng trở lên trong tổng số $48$ trận đấu là $p = 0,0002$ (tương đương tỷ lệ $1$ trên $5.000$). Kết quả cho thấy sự vượt trội mà dung hợp mang lại.
\subsection{Ảo giác mô hình và tính xác thực của kết quả}

Bảng~\ref{tab:axisB3} báo cáo số liệu thống kê về tỷ lệ chấp thuận logic ở cấp độ mệnh đề trên cả năm phương án đánh giá. Kết quả dung hợp đạt tỷ lệ chấp thuận là $94,4\%$, minh chứng rằng quy trình dung hợp bảo toàn tốt tính trung thực của dữ kiện trong khi vẫn thực hiện tổng hợp thông tin từ nhiều bản thảo đề xuất khác nhau.

\begin{table}[H]
  \centering
  \caption{Trục B.3 --- đánh giá tính xác thực qua quan hệ kéo theo cấp mệnh đề (sử dụng mô hình phán quyết \texttt{gpt-4o-mini}).}
  \label{tab:axisB3}
  \begin{tabularx}{\linewidth}{X r C C C}
    \toprule
    Mô hình & $n$ & tỷ lệ chấp thuận \% & Ảo giác TB & Trường hợp xấu nhất \\
    \midrule
    \textbf{MoA FUSION}            & \textbf{154} & \textbf{94,4} & - & - \\
    claude-haiku-4.5             & 100 & 95,5 & 0,07 & 1 \\
    gemini-2.5-flash             & 100 & 94,7 & 0,10 & 2 \\
    gpt-4o-mini                  & 100 & 92,8 & 0,08 & 2 \\
    gpt-4o          & 50  & 89,9 & 0,08 & 1 \\
    \bottomrule
  \end{tabularx}
\end{table}

Trên toàn bộ cơ sở dữ liệu gồm $923$ mệnh đề được đánh giá, tỷ lệ chấp thuận tổng hợp đạt $94,4\%$. Tỷ lệ chấp thuận thấp nhất ghi nhận ở mô hình GPT-4o đơn lẻ ($89,9\%$). Kết quả này củng cố quan sát thực tế rằng khi hoạt động độc lập không có bước ``kiểm duyệt biên tập'' mà pipeline MoA cung cấp một cách gián tiếp, GPT-4o có xu hướng vận hành ở chế độ sáng tạo hơn và đôi khi đưa vào những dữ kiện không được văn bản gốc hỗ trợ trực tiếp, hay còn gọi là ảo giác.

\section{Kết quả Trục C --- Đánh giá của con người}

\subsection{Thiết lập khảo sát}

Một khảo sát đã được triển khai với sự tham gia của $17$ người đánh giá độc lập, bao phủ toàn bộ $48$ tác vụ thực nghiệm, thu về tổng cộng $192$ kết quả xếp hạng mù $K = 3$. Mỗi tác vụ hiển thị đồng thời ba ứng viên tóm tắt --- bản dung hợp MoA, bản GPT-4o đơn lẻ, và bản thảo đề xuất GPT-4o-mini --- dưới các nhãn ngẫu nhiên \textit{Bản A}, \textit{Bản B}, \textit{Bản C} nhằm ẩn hoàn toàn danh tính của mô hình.

\subsection{Kết quả nổi bật trên Trục C}

\begin{table}[H]
  \centering
  \caption{Trục C --- các chỉ số thống kê mô tả tổng hợp từ người đánh giá, quy mô $n = 192$ kết quả xếp hạng trên $48$ tác vụ.}
  \label{tab:axisC}
  \begin{tabularx}{\linewidth}{X C C C}
    \toprule
    Mô hình & Hạng trung bình ($\downarrow$) & Tỷ lệ xếp hạng nhất & Tỷ lệ đối đầu vs. Dung hợp \\
    \midrule
    \textbf{MoA FUSION}    & \textbf{1,95} & \textbf{40,9\%} & --- \\
    gpt-4o     & 2,00 & 32,3\% & Dung hợp thắng 40,9\% \\
    gpt-4o-mini & 2,05 & 26,8\% & Dung hợp thắng 40,9\% \\
    \bottomrule
  \end{tabularx}
\end{table}

Bản tóm tắt của MoA FUSION là phương án được người đọc ưa thích nhất trên tất cả các tiêu chí đo lường tổng hợp: thứ hạng trung bình thấp nhất, tỷ lệ xếp hạng nhất cao nhất và ưu thế chiến thắng trong các cặp đối đầu so sánh trực diện. Tuy nhiên, khoảng cách vượt trội này nhỏ hơn rất nhiều so với những gì Trục B thể hiện. LLM-Judge báo cáo tỷ lệ thắng lên tới $68,9\%$ trong các phán quyết so sánh quyết định (dung hợp đối đầu bản thảo tốt nhất); trong khi kết quả so sánh của con người chỉ báo cáo mức độ ưa thích $40,9\%$ dành cho bản dung hợp. Khoảng cách chênh lệch này đã được nghiên cứu, tiếp tục thảo luận và lý giải ở các mục tiếp theo.

\subsection{Tương tác theo nhóm chủ đề}

Bảng~\ref{tab:axisC-topic} trình bày kết quả phân tích chi tiết theo từng nhóm chủ đề. Để phục vụ phân tích sâu, năm chủ đề báo chí của \textit{tienphong.vn} được gom lại thành ba nhóm thể loại dựa trên nhận định trực quan của người đọc: nhóm chủ đề HÀNH CHÍNH (chính trị, hành chính, công văn nghị quyết; gồm $21$ bài), nhóm chủ đề GIẢI TRÍ (lễ hội, người nổi tiếng, giải trí; gồm $13$ bài), và nhóm chủ đề HỖN HỢP (hình sự, tin quốc tế, quan điểm độc giả; gồm $14$ bài).

\begin{table}[H]
  \centering
  \caption{Trục C --- sở thích xếp hạng phân nhánh theo nhóm chủ đề. Ký hiệu: $F$ = Bản dung hợp, $4o$ = GPT-4o đơn lẻ, $m$ = Bản thảo gpt-4o-mini.}
  \label{tab:axisC-topic}
  \begin{tabularx}{\linewidth}{X r C C C C C}
    \toprule
    Nhóm chủ đề & $n$ & Tỷ lệ $F$ hạng nhất & Tỷ lệ $4o$ hạng nhất & Tỷ lệ $m$ hạng nhất & $F > 4o$ & $F > m$ \\
    \midrule
    HÀNH CHÍNH     & 84 & 23\% & 37\% & \textbf{41\%} & 36\% & 43\% \\
    GIẢI TRÍ       & 52 & \textbf{58\%} & 21\% & 21\% & 65\% & 65\% \\
    HỖN HỢP        & 56 & \textbf{59\%} & 29\% & 13\% & 70\% & 77\% \\
    \bottomrule
  \end{tabularx}
\end{table}

Trên nhóm nội dung mang tính hành chính (đại hội Đảng, báo cáo hành chính, hội nghị ban ngành), người đọc có xu hướng ưa thích bản tóm tắt ngắn nhất (bản thảo tạo bởi GPT-4o-mini) --- vốn chính là ứng viên bị LLM-Judge xếp cuối cùng về tiêu chí Độ bao phủ (Coverage) trong Mục~4.3.1. Bản tóm tắt dung hợp --- một văn bản dài hơn, được tổng hợp biên tập chi tiết --- bị người đọc đánh giá thấp hơn; do thiếu đi sự hứng thú tự thân với chủ đề. Ngược lại, đối với thể loại giải trí và các chủ đề hỗn hợp, bản tóm tắt dung hợp đầy đủ thông tin hơn lại được ưa thích vượt trội (đạt tỷ lệ xếp hạng nhất lần lượt là $58\%$ và $59\%$), hoàn toàn đồng điệu với nhận định của LLM-Judge.

LLM-Judge hoàn toàn không có khả năng nhận biết mức độ hấp dẫn của chủ đề đối với người đọc thực tế; nó chỉ chấm điểm máy móc theo rubric định sẵn và luôn thưởng điểm cho Độ bao phủ một cách đồng nhất. Trong khi đó, con người mang theo mục đích đọc thực tế khi tiếp cận bài viết, và mục đích này sẽ triệt tiêu tiêu chí Độ bao phủ khi họ gặp các chủ đề tẻ nhạt. Đây chính là phát hiện về mặt phương pháp luận mà khung đánh giá ba trục được thiết kế để chỉ ra.

\subsection{Độ đồng thuận giữa các người đánh giá}

Tính trung bình trên $48$ tác vụ, số lượng người đánh giá đồng thuận lớn nhất về vị trí hạng nhất (\#1) đạt $2{,}69$ trên $4$ người ($\approx 67\%$). Có 8 tác vụ đạt sự đồng thuận tuyệt đối (4/4), $17$ tác vụ ghi nhận sự đồng thuận của 3 trên 4 người, $23$ tác vụ bị chia rẽ nửa này nửa kia (2 vs. 2), và hoàn toàn không có tác vụ nào có ý kiến bị phân mảnh hoàn toàn. Kết quả này tương ứng với hệ số Fleiss' $\kappa$ ước tính nằm trong khoảng đồng thuận vừa phải ($\approx 0{,}30$--$0{,}40$); mức độ đồng thuận thấp nhất ở nhóm chủ đề HÀNH CHÍNH (trung bình chỉ $2{,}48$ người đồng ý với nhau) và cao nhất ở nhóm GIẢI TRÍ ($2{,}92$ người). Chỉ số này phản ánh chính xác hành vi của độc giả trong thực tế --- họ chia sẻ một số sở thích chung cốt lõi, nhưng cũng có những bất đồng sâu sắc khi đánh giá các nội dung hành chính khô khan.

\section{Phân tích đa trục}

Góc nhìn khoa học quan trọng nhất từ các thực nghiệm chính là việc đối chiếu đồng thời cả ba trục đánh giá. Bảng~\ref{tab:cross-axis} đối chiếu phép so sánh then chốt giữa bản dung hợp MoA và GPT-4o chạy đơn lẻ trên từng trục.

\begin{table}[H]
  \centering
  \caption{Đối chiếu đa trục, bản dung hợp đối đầu với các mô hình cơ sở.}
  \label{tab:cross-axis}
  \begin{tabularx}{\linewidth}{l X X}
    \toprule
    Trục & Chỉ số đo lường & Kết quả phán quyết \\
    \midrule
    A   & ROUGE-1 ($0{,}4240$ đấu với $0{,}3621$)   & Dung hợp dẫn $+0{,}062$ \\
    A   & BERTScore ($0{,}6759$ đấu với $0{,}6246$) & Dung hợp dẫn $+0{,}051$ \\
    B.1 & Điểm Tổng thể Rubric ($4{,}96$ đấu với $4{,}88$) & Dung hợp dẫn $+0{,}08$ \\
    B.2a & Tỷ lệ thắng đấu cặp so với bản thảo tốt nhất ($68,9\%$) & Dung hợp thắng vượt trội \\
    C   & Kết quả so sánh của con người ($40,9\%$ đa số tương đối) & Dung hợp thắng với chênh lệch thấp \\
    \bottomrule
  \end{tabularx}
\end{table}

Cả ba góc nhìn độc lập đều thống nhất rằng phương pháp dung hợp MoA mang lại kết quả tốt hơn. Sự đồng thuận về mặt chiều hướng này là kết quả lý tưởng nhất của khung đánh giá ba trục.

Biên độ vượt trội của phương pháp dung hợp biến động rất lớn tùy theo hệ số đo lường: tăng nhẹ $+0{,}05$ đến $+0{,}06$ trên Trục A, dẫn $+0{,}08$ trên Trục B.1, thắng áp đảo với cách biệt $+18{,}9$ điểm phần trăm trên Trục B.2a (vượt mức $50\%$), nhưng lại thắng với chênh lệch nhỏ trên Trục C ($40,9\%$ chưa đạt mức đa số tuyệt đối, dù vẫn chiếm đa số tương đối). Khoảng chênh lệch $28$ điểm phần trăm giữa Trục B.2a ($68,9\%$) và Trục C ($40,9\%$) là số liệu đáng giá nhất trong khóa luận này: \textbf{Sự ưu tiên của LLM-Judge, trong thiết lập thực nghiệm này, không phản ánh đúng biên độ ưa thích của con người thực tế}. Chiều hướng đánh giá là chính xác, nhưng biên độ đã bị phóng đại quá mức.

Khoảng chênh lệch này có thể được giải thích bằng các giả thuyết thu thập được sau quá trình phỏng vấn người đánh giá: 
\begin{itemize}
  \item Thiên kiến ủng hộ độ dài và tính toàn diện: LLM-Judge luôn cộng điểm cho Độ bao phủ (Coverage) một cách nhất quán; bản tóm tắt dung hợp dài hơn và bao hàm nhiều thông tin hơn bản GPT-4o đơn lẻ, nên LLM-Judge dễ dàng phán quyết. Con người, khi đối diện với các chủ đề hành chính khô khan, lại phản ứng ngược lại.
  \item Tương tác theo nhóm chủ đề: LLM-Judge không có nhận thức về sở thích và hành vi đọc; con người thì có (Mục~4.4.3).
  \item Hiệu ứng phân cực phán quyết:  Các LLM-Judge khi so sánh thường nhạy cảm với các khác biệt nhỏ và có xu hướng phóng đại khoảng cách này \cite{Zheng2023}; trong khi con người thực tế đánh giá một cách dè dặt và có xu hướng nghiêng về kết quả hòa nhiều hơn.
\end{itemize}

Việc chỉ sử dụng một trục đánh giá duy nhất là không đủ để báo cáo đầy đủ về hiệu năng của phương pháp MoA. Kết quả từ cả ba trục cần được kết hợp chặt chẽ để bổ trợ lẫn nhau; nếu chỉ chạy riêng Trục B, khóa luận sẽ vội vã kết luận tỷ lệ thắng áp đảo lên tới $77\%$ trong khi thực tế kết quả ghi nhận từ con người (Trục C) chỉ dừng lại ở mức $41\%$.

\section{Hạn chế phương pháp luận}

Sau quá trình thực nghiệm và đánh giá toàn diện, có một số hạn chế phương pháp luận có thể ghi nhận như sau: 

\begin{itemize}
    \item Kiến trúc đơn giản hoá: Cấu hình MoA được triển khai thực tế trong hệ thống này là cấu hình tối thiểu mà nghiên cứu gốc xem xét: chỉ gồm hai tầng gồm ba mô hình đề xuất kết hợp một mô hình tổng hợp. Bảng~3 và 4 trong công trình của Wang et al.\ \cite{Wang2024} chỉ ra rằng hiệu năng của hệ thống tăng đơn điệu theo độ sâu của tầng tác nhân ($\ell \in \{1, 2, 3\}$) và số lượng mô hình đề xuất trong mỗi tầng ($n \in \{3, 4, 5, 6\}$). Thiết kế hiện tại được lựa chọn để tối ưu hóa chi phí và kiểm soát độ trễ thời gian thực; việc nâng cấp chiều sâu và chiều rộng của mạng lưới tác nhân là những hướng phát triển tiềm năng cho các công trình tương lai.
    \item Quy mô nghiên cứu con người và động lực của người đánh giá:
Nghiên cứu con người hiện tại có sự tham gia của $17$ người đánh giá độc lập, cung cấp một cơ sở định tính rất giá trị. Tuy nhiên, việc phụ thuộc vào nhóm người đánh giá thực tế này cũng đặt ra một số thách thức đặc thù. Thứ nhất là vấn đề sở thích chủ đề cá nhân; khi đối mặt với các bài báo, những người đánh giá ngẫu nhiên có thể thể hiện mức độ tập trung nhận thức rất khác nhau tùy thuộc vào mức độ quan tâm của họ đối với nội dung bài viết. Sự thiên lệch về hứng thú này có thể đưa thêm các yếu tố nhiễu vào số liệu thống kê sở thích chung.
Thứ hai, nghiên cứu cần kiểm soát hiện tượng mệt mỏi nhận thức và sự suy giảm năng lượng của người đánh giá theo thời gian. Khi các phiên đánh giá kéo dài, tài nguyên chú ý của người tham gia tự động giảm đi, dẫn đến sự thiếu nhất quán trong việc xếp hạng hoặc xu hướng đưa ra các phán quyết an toàn (như chọn hòa) ở các tác vụ cuối cùng. Mặc dù mẫu dữ liệu hiện tại đủ lớn để khẳng định các xu hướng chính, các nghiên cứu tiếp theo sẽ đạt chất lượng cao hơn nếu thực hiện trong môi trường kiểm soát chặt chẽ hơn để giảm thiểu các biến động hành vi này và đảm bảo tính đồng đều trong đánh giá trên toàn bộ tập dữ liệu.

\item Hệ số overlap tính trực tiếp với văn bản nguồn: Các hệ số ROUGE / BLEU / BERTScore được tính toán so với bài viết gốc, không phải so với bản tóm tắt chuẩn do con người viết, vì hiện tại chưa có bộ dữ liệu benchmark tóm tắt tiếng Việt nào được chú thích bởi con người ở quy mô mà dự án này yêu cầu. Đây là thực tế chung trong các nghiên cứu NLP học thuật tại Việt Nam và đã được khoá luận này thừa nhận như một hạn chế về mặt phương pháp luận của Trục A. Việc đối chiếu đa trục thực hiện trong Mục~4.5 là chốt chặn thực nghiệm quan trọng nhằm tránh việc diễn giải quá đà các điểm số overlap này.
\end{itemize}

\section{Tổng kết chương}

Chương này đã hoàn thành việc kiểm chứng thực nghiệm toàn diện hệ thống trên 154 bài báo tiếng Việt cân bằng theo năm nhóm chủ đề, và thu được hai
nhóm kết quả chính. Nhóm thứ nhất là sự đồng thuận về chiều hướng: cả ba trục đánh giá đều thống nhất rằng phương pháp dung hợp MoA tạo ra bản tóm tắt chất lượng cao hơn bất kỳ mô hình đơn lẻ nào, với ROUGE-1 dẫn $+0{,}062$, BERTScore dẫn $+0{,}051$, điểm Tổng thể Rubric dẫn $+0{,}08$, tỷ lệ thắng đấu cặp LLM-Judge đạt $68{,}9\%$, và tỷ lệ ưu tiên của con người đạt $40{,}9\%$ đa số tương đối.
 Đặc biệt, trong phép so sánh có kiểm soát mô hình (dung hợp MoA đối đầu GPT-4o đơn lẻ, cùng là đầu ra của GPT-4o), bản dung hợp thắng $37/48$ lần ($p = 0{,}0002$), chứng minh rằng hiệu quả đến từ chính quá trình dung hợp chứ không phải từ sự chênh lệch năng lực mô hình. 
 Nhóm thứ hai là các phát hiện phương pháp luận: khoảng cách 28 điểm phần trăm giữa tỷ lệ thắng của LLM-Judge ($68{,}9\%$) và tỷ lệ ưu tiên của con người ($40{,}9\%$) cho thấy LLM-Judge phóng đại biên độ vượt trội; đồng thời, hiện tượng đảo ngược sở thích theo nhóm chủ đề --- con người ưa thích bản tóm tắt ngắn nhất cho các bài hành chính khô khan dù LLM-Judge luôn thưởng cho độ bao phủ --- khẳng định rằng không một trục đánh giá đơn lẻ nào đủ để phản ánh đầy đủ chất lượng tóm tắt. 
 Chương~5 sẽ tổng hợp các đóng góp khoa học, tái khẳng định các hạn chế và đề xuất các hướng phát triển tương lai.
